% NOTE (7 Oct 2026): the numbers in this draft are CLIP-LEVEL. The thesis now reports
% film-level macro-F1 as the headline (results/film_level.json; docs/current_state.md,
% section 10). Use the film-level numbers when writing the final version.
% =============================================================================
% HELD BACK from the thesis on 6 Oct 2026 (to be rewritten by the author).
% Not part of the Overleaf build. Was the end of Section "Aim of this thesis" in chapter/0_intro.tex.
% =============================================================================

To answer these questions, this thesis makes the following contributions.

First, it builds a \emph{two-stage, interpretable pipeline} from audio to genre. A frozen
pretrained audio network turns each clip into an embedding; a regressor trained on
listener ratings maps the embedding to eight emotions (valence, energy, tension, anger,
fear, happiness, sadness and tenderness); and a linear classifier maps those emotions,
with three simple combinations of them, to a multi-label genre vector. Every decision of
the genre classifier can therefore be read in terms of named emotions.

Second, it compares the pipeline against a \emph{direct-audio baseline built from five
representations rather than one}: three general-purpose learned embeddings (VGGish, the
Audio Spectrogram Transformer and CLAP), one embedding pretrained specifically on music
(MusiCNN), and a set of hand-crafted music-information-retrieval descriptors, with a raw
waveform model (wav2vec~2.0) as a sixth point of reference. A result against a single
baseline invites the objection that the baseline was badly chosen; a result against five
representations of different architecture and training data does not.

Third, it includes a \emph{control} that separates the effect of emotion from the effect
of compression. The emotion route reduces a high-dimensional embedding to a handful of
numbers, and with little training data fewer numbers can help on their own. The control
compresses the same embedding to eight principal components, which are just as compact
but mean nothing. If emotion helps only by being small, the control helps equally.

Fourth, it tests the pipeline \emph{across datasets}. Besides the film-music dataset of
\citet{eerola2011comparison}, on which the emotion model is trained, it uses the
independent dataset of blockbuster films published by \citet{ma2021computational}, which
comes without emotion ratings and without audio. The two datasets are mapped to a shared
genre space, the published in-domain result on the second dataset is reproduced first, and
models are then trained on one dataset and evaluated on the other in both directions.
This asks whether what the models learn is a property of film music or of one collection
of films.

Finally, the evaluation itself is designed so that its numbers can be trusted: every split
is grouped by film, every score is paired with a chance floor from its own label space,
every in-domain comparison is reported under two significance tests that answer
different questions and every transfer comparison under a film bootstrap, together with a check of how far the conclusions depend on the choice of
metric. Where this apparatus weakened an earlier result, the weaker result is the one
reported.
